Knowledge tracing (KT) assumes that every assessment opportunity is linked to
one or more knowledge components (KCs), yet grammar curricula commonly express
learning targets as prose descriptors rather than operational skills.  This
report presents \system, an auditable framework for constructing and comparing
grammar KC representations.  A frozen pilot design selects 139 descriptors
from a digest-declared English Grammar Profile snapshot and uses constrained,
two-phase model-assisted normalisation to map supported claims into a closed
six-dimensional representation of English verbal morphosyntax.  Deterministic
realisation then constructs a KC-independent item bank.  A separate structural
oracle generates one fixed synthetic event stream before any evaluated KC
policy is selected, so alternative representations are compared on identical
items and outcomes.  Development-only candidate diagnostics, frozen KC
policies, Q-matrix construction, and non-updating compositional probes make
support, coverage, and transfer assumptions explicit.  The present revision
documents the implemented methodology and its evidence requirements;
quantitative research comparisons are deliberately reserved for a fresh,
validated run of the final pipeline.  The contribution is therefore a
controlled experimental design for studying grammar representations, not a
claim that synthetic performance establishes the cognitive reality of any KC.
